\documentclass[12pt, letterpaper]{article}
\usepackage[T1]{fontenc}
\usepackage{lmodern}
\usepackage[margin=0.5in]{geometry} 
\renewcommand{\familydefault}{\sfdefault}
\usepackage{amsmath}




\begin{document}

\begin{center}
    {\Large\bfseries Assignment 18: Factoring Quadratics}\\[4pt]
    {\large\bfseries DO NOT WRITE ON THIS PAPER}\\[4pt]
\end{center}



\fbox{%
  \begin{minipage}[t][0.40\textheight][t]{.95\textwidth}
    \begin{center}
        \textbf{Part A}
    \end{center}
    
    Use the following equation to fill out the t-table below.

    \begin{center}
        \(y = x^2 + 8x + 12\)
    \end{center}

    \begin{center}
        \renewcommand{\arraystretch}{1.5}
        \begin{tabular}{|p{1cm}|p{1cm}|}
        \hline
        $x$ & $y$ \\ \hline
        $-6$ &  \\ \hline
        $-5$ &  \\ \hline
        $-4$ &  \\ \hline
        $-3$ &  \\ \hline
        $-2$ &  \\ \hline
    \end{tabular}
    \end{center}
    \end{minipage}%
}%


\fbox{%
  \begin{minipage}[t][0.46\textheight][t]{.95\textwidth}
    \begin{center}
        \textbf{Part B}
    \end{center}

    Solve the quadratic equation \( x^2 + 8x + 12 = 0 \) and show your work.

    \vspace{1em}

    Do you recognize any values from the T-table in part A? If so what values and what are special about them? 
    \end{minipage}%
}%


\newpage

\fbox{%
  \begin{minipage}[t][0.40\textheight][t]{.95\textwidth}
    \begin{center}
        \textbf{Part C}
    \end{center}
    
   Factor the following quadratic expressions and show your work.

   \vspace{1em}

   \begin{enumerate}
       \item Factor by grouping: \( 5y^2 + 12y - 9 \)
       \item Difference of Squares: \( 49w^2 - 4 \)
    \end{enumerate}
    \end{minipage}%
}

\fbox{%
  \begin{minipage}[t][0.46\textheight][t]{.95\textwidth}
    \begin{center}
        \textbf{Question 1} 
    \end{center}



    \textbf{Big Question:} What does it mean to factor a quadratic completely? Why do we try and factor to solve quadratic equations?

        \vspace{1em}

       Use the phrases ``Zero-Product Property, Linear Terms, and Product'' in your answer.


       \vspace{1em}
       Then factor the following quadratic expression completely and show your work: \( 3x^2 -2x - 21 \)

    \vfill
    \hfill
  \end{minipage}%
}%

\newpage

\fbox{%
  \begin{minipage}[t][0.42\textheight][t]{.7\textwidth}
    \begin{center}
        \textbf{Question 2 Work}
    \end{center}

    \begin{center}
    If I factor a quadratic expression how can I check my work? 
    
    \vspace{1em}
    If I solve a quadratic equation how can I check my work?

    \vspace{1em}

    \emph{Hint: Think about the example from the slides: $x^2-3x+2$.}

    \end{center}
    \hfill
    \vfill
  \end{minipage}%
}%
\fbox{%
  \begin{minipage}[t][0.42\textheight][t]{0.24\textwidth}
    \begin{center}
      {\bfseries Question 2 Answer}
      % Question 2 answer box -  no need to edit
        \vspace*{\fill}
      
      {\Large Write your answer in your Class Notes}
    \end{center}
    \vspace*{\fill} 
  \end{minipage}%
}

\fbox{%
  \begin{minipage}[t][0.42\textheight][t]{.7\textwidth}
    \begin{center}
        \textbf{Question 3 Work}
    \end{center}

    \begin{center}
    Solve \(2x^2 + x -21 = 0\).
    
    \vspace{1em}
    Show your work and check your answer.
    \end{center}

  \end{minipage}%
}%
\fbox{%
  \begin{minipage}[t][0.42\textheight][t]{0.24\textwidth}
    \begin{center}
      {\bfseries Question 3 Answer}
      
      % Question 3 answer box -  no need to edit
        \vspace*{\fill}
      
      {\Large Write your answer in your Class Notes}
    \end{center}
    \vspace*{\fill}
  \end{minipage}%
}

\newpage

\begin{center}
    \textbf{Additional Space}
\end{center}

\vspace{0.4em}

\noindent\makebox[\textwidth][c]{%
\fbox{%
  \rule{0pt}{0.20\textheight}\rule{7.35in}{0pt}%
}%
}

\vspace{0.45em}

\noindent\makebox[\textwidth][c]{%
\fbox{%
  \rule{0pt}{0.20\textheight}\rule{7.35in}{0pt}%
}%
}

\vspace{0.45em}

\noindent\makebox[\textwidth][c]{%
\fbox{%
  \rule{0pt}{0.20\textheight}\rule{7.35in}{0pt}%
}%
}

\vspace{0.45em}

\noindent\makebox[\textwidth][c]{%
\fbox{%
  \rule{0pt}{0.20\textheight}\rule{7.35in}{0pt}%
}%
}

\end{document}
